\documentclass[10pt,a4paper]{amsart}
\usepackage{amsmath,amssymb,amsthm,graphicx}
\usepackage{verbatim}
\newenvironment{reference}{\comment}{\endcomment}
%\newenvironment{reference}{}{}
\newenvironment{slogan}{\comment}{\endcomment}
\newenvironment{history}{\comment}{\endcomment}

% For commutative diagrams we use Xy-pic
\usepackage[all]{xy}

% We use 2cell for 2-commutative diagrams.
\xyoption{2cell}
% Theorem environments.
%
\theoremstyle{plain}
\newtheorem{theorem}[subsection]{Theorem}
\newtheorem{proposition}[subsection]{Proposition}
\newtheorem{lemma}[subsection]{Lemma}

\theoremstyle{definition}
\newtheorem{definition}[subsection]{Definition}
\newtheorem{example}[subsection]{Example}
\newtheorem{exercise}[subsection]{Exercise}
\newtheorem{situation}[subsection]{Situation}

\theoremstyle{remark}
\newtheorem{remark}[subsection]{Remark}
\newtheorem{remarks}[subsection]{Remarks}

\numberwithin{equation}{subsection}

% Macros
%
\def\lim{\mathop{\mathrm{lim}}\nolimits}
\def\colim{\mathop{\mathrm{colim}}\nolimits}
\def\Spec{\mathop{\mathrm{Spec}}}
\def\Hom{\mathop{\mathrm{Hom}}\nolimits}
\def\Ext{\mathop{\mathrm{Ext}}\nolimits}
\def\SheafHom{\mathop{\mathcal{H}\!\mathit{om}}\nolimits}
\def\SheafExt{\mathop{\mathcal{E}\!\mathit{xt}}\nolimits}
\def\Sch{\mathit{Sch}}
\def\Mor{\mathop{\mathrm{Mor}}\nolimits}
\def\Ob{\mathop{\mathrm{Ob}}\nolimits}
\def\Sh{\mathop{\mathit{Sh}}\nolimits}
\def\NL{\mathop{N\!L}\nolimits}
\def\CH{\mathop{\mathrm{CH}}\nolimits}
\def\proetale{{pro\text{-}\acute{e}tale}}
\def\etale{{\acute{e}tale}}
\def\QCoh{\mathit{QCoh}}
\def\Ker{\mathop{\mathrm{Ker}}}
\def\Im{\mathop{\mathrm{Im}}}
\def\Coker{\mathop{\mathrm{Coker}}}
\def\Coim{\mathop{\mathrm{Coim}}}

% Boxtimes
%
\DeclareMathSymbol{\boxtimes}{\mathbin}{AMSa}{"02}

%
% Macros for moduli stacks/spaces
%
\def\QCohstack{\mathcal{QC}\!\mathit{oh}}
\def\Cohstack{\mathcal{C}\!\mathit{oh}}
\def\Spacesstack{\mathcal{S}\!\mathit{paces}}
\def\Quotfunctor{\mathrm{Quot}}
\def\Hilbfunctor{\mathrm{Hilb}}
\def\Curvesstack{\mathcal{C}\!\mathit{urves}}
\def\Polarizedstack{\mathcal{P}\!\mathit{olarized}}
\def\Complexesstack{\mathcal{C}\!\mathit{omplexes}}
% \Pic is the operator that assigns to X its picard group, usage \Pic(X)
% \Picardstack_{X/B} denotes the Picard stack of X over B
% \Picardfunctor_{X/B} denotes the Picard functor of X over B
\def\Pic{\mathop{\mathrm{Pic}}\nolimits}
\def\Picardstack{\mathcal{P}\!\mathit{ic}}
\def\Picardfunctor{\mathrm{Pic}}
\def\Deformationcategory{\mathcal{D}\!\mathit{ef}}
\usepackage{cleveref}
\crefname{originalequation}{original equation}{original equations}
\Crefname{originalequation}{Original equation}{Original equations}
\makeatletter
\newlabel{equation-sheafy}{{0A9L}{source}{}{sourcepointer-1}{}}
\newlabel{equation-sheafy@cref}{{[originalequation][0][]0A9L}{[1][0][]source}{}{}{}}
\makeatother
\crefname{originallemma}{original lemma}{original lemmas}
\Crefname{originallemma}{Original lemma}{Original lemmas}
\makeatletter
\newlabel{perfect-lemma-quasi-coherence-direct-image}{{08D5}{source}{}{sourcepointer-2}{}}
\newlabel{perfect-lemma-quasi-coherence-direct-image@cref}{{[originallemma][0][]08D5}{[1][0][]source}{}{}{}}
\makeatother
\crefname{originaldefinition}{original definition}{original definitions}
\Crefname{originaldefinition}{Original definition}{Original definitions}
\makeatletter
\newlabel{perfect-definition-supported-on}{{08DA}{source}{}{sourcepointer-3}{}}
\newlabel{perfect-definition-supported-on@cref}{{[originaldefinition][0][]08DA}{[1][0][]source}{}{}{}}
\makeatother
\crefname{originallemma}{original lemma}{original lemmas}
\Crefname{originallemma}{Original lemma}{Original lemmas}
\makeatletter
\newlabel{perfect-lemma-generator-with-support}{{0A9A}{source}{}{sourcepointer-4}{}}
\newlabel{perfect-lemma-generator-with-support@cref}{{[originallemma][0][]0A9A}{[1][0][]source}{}{}{}}
\makeatother
\crefname{originallemma}{original lemma}{original lemmas}
\Crefname{originallemma}{Original lemma}{Original lemmas}
\makeatletter
\newlabel{derived-lemma-write-as-colimit}{{09SN}{source}{}{sourcepointer-5}{}}
\newlabel{derived-lemma-write-as-colimit@cref}{{[originallemma][0][]09SN}{[1][0][]source}{}{}{}}
\makeatother
\title{The sheafyness condition is equivalent to the right adjoint carrying supported objects to supported objects}
\author{The Stacks Project Authors}
\date{stacks-project:duality.tex:local-snapshot-20261012; GNU FDL 1.2 or later}
\begin{document}
\maketitle
\noindent\textit{Editorial note.} All delivered mathematical text is the author's own, copied from the pinned Stacks Project chapter file duality.tex; no statement or proof step was rewritten, reordered or removed.\par\medskip
\noindent\textit{Editorial note.} Preamble provenance: the pinned chapter file opens with an input of the shared preamble and carries no preamble text of its own, so the project's single Stacks-wide preamble staged at sources/stacks-duality/preamble.tex is excerpted verbatim under the recipe fixed by LANE\_BRIEF 22.1, as three spans, lines 8-13 (verbatim package and the reference/slogan/history environments), lines 14-18 (Xy-pic setup) and lines 154-217 (the native theorem declarations and mathematical macros). Nothing was added, removed or redefined; the wrapper class and the amsmath/amssymb/amsthm/graphicx base packages are supplied by the assembler.\par\medskip
\noindent\textit{Editorial note.} The counted statement and proof contain no cross reference, no citation and no included file, so the delivered body is self contained.\par\medskip
\noindent\textit{Editorial note.} Canonical location of the counted proposition: Stacks tag 0A9N (sources/stacks-duality/pinned-original-source.tex line 1035).\par\medskip
\noindent\textit{Allowed standard prerequisites (editorial).}\begin{itemize}
\item The right adjoint of pushforward and the sheafyness condition. Same chapter, equation labelled equation-sheafy, invoked in the statement, with the right adjoint a of Rf\_* as in the statement. Fixes the condition being characterised.
\item Quasi-coherence of the direct image and the notion of being supported on a closed subset. Derived Categories of Schemes, Lemma labelled perfect-lemma-quasi-coherence-direct-image and Definition labelled perfect-definition-supported-on, cited at proof lines 1058-1062. Make the third term of the triangle quasi-coherent and supported.
\item Generation of the complexes supported on a closed subset by a perfect object. Derived Categories of Schemes, Lemma labelled perfect-lemma-generator-with-support, cited at proof lines 1088-1092. Provides E and the generating property used in (3).
\item Every object as a derived colimit of a system with cones direct sums of shifts. Derived Categories, Lemma labelled derived-lemma-write-as-colimit, cited at proof lines 1092-1096. Feeds the induction that the subcategory is everything.
\end{itemize}\medskip
\ifdefined\hypertarget\hypertarget{sourcepointer-1}{}\fi
\noindent\textit{Source reference.} Original equation 0A9L: sources/stacks-duality/pinned-original-source.tex, source lines 988--988. Same-chapter pointer into the author environment equation at source line 988; invoked at proof line 1072 of the counted proof. It is retained only as a navigation pointer; every load-bearing step is inside the delivered span.\par\medskip
\ifdefined\hypertarget\hypertarget{sourcepointer-2}{}\fi
\noindent\textit{Source reference.} Original lemma 08D5: sources/stacks-chow/tags.txt, source lines 10625--10625. Cross-chapter standard prerequisite located in the project tag list staged in the local snapshot; invoked at proof line 1062 of the counted proof. It is retained only as a navigation pointer; every load-bearing step is inside the delivered span.\par\medskip
\ifdefined\hypertarget\hypertarget{sourcepointer-3}{}\fi
\noindent\textit{Source reference.} Original definition 08DA: sources/stacks-chow/tags.txt, source lines 10630--10630. Cross-chapter standard prerequisite located in the project tag list staged in the local snapshot; invoked at proof line 1065 of the counted proof. It is retained only as a navigation pointer; every load-bearing step is inside the delivered span.\par\medskip
\ifdefined\hypertarget\hypertarget{sourcepointer-4}{}\fi
\noindent\textit{Source reference.} Original lemma 0A9A: sources/stacks-chow/tags.txt, source lines 12943--12943. Cross-chapter standard prerequisite located in the project tag list staged in the local snapshot; invoked at proof line 1094 of the counted proof. It is retained only as a navigation pointer; every load-bearing step is inside the delivered span.\par\medskip
\ifdefined\hypertarget\hypertarget{sourcepointer-5}{}\fi
\noindent\textit{Source reference.} Original lemma 09SN: sources/stacks-chow/tags.txt, source lines 12361--12361. Cross-chapter standard prerequisite located in the project tag list staged in the local snapshot; invoked at proof line 1096 of the counted proof. It is retained only as a navigation pointer; every load-bearing step is inside the delivered span.\par\medskip
\begin{lemma}
\label{lemma-when-sheafy}
Let $f : X \to Y$ be a morphism of quasi-compact and quasi-separated
schemes. Let $a$ be the right adjoint to
$Rf_* : D_\QCoh(\mathcal{O}_X) \to D_\QCoh(\mathcal{O}_Y)$.
Let $V \subset Y$ be quasi-compact open with inverse image $U \subset X$.
\begin{enumerate}
\item For every $Q \in D_\QCoh^+(\mathcal{O}_Y)$
supported on $Y \setminus V$ the image $a(Q)$ is supported on
$X \setminus U$ if and only if (\ref{equation-sheafy})
is an isomorphism on all $K$ in $D_\QCoh^+(\mathcal{O}_Y)$.
\item For every $Q \in D_\QCoh(\mathcal{O}_Y)$
supported on $Y \setminus V$ the image $a(Q)$ is supported on
$X \setminus U$ if and only if (\ref{equation-sheafy})
is an isomorphism on all $K$ in $D_\QCoh(\mathcal{O}_Y)$.
\item If $a$ commutes with direct sums, then the equivalent conditions of
(1) imply the equivalent conditions of (2).
\end{enumerate}
\end{lemma}

\begin{proof}
Proof of (1). Let $K \in D_\QCoh^+(\mathcal{O}_Y)$.
Choose a distinguished triangle
$$
K \to Rj_*K|_V \to Q \to K[1]
$$
Observe that $Q$ is in $D_\QCoh^+(\mathcal{O}_Y)$
(Derived Categories of Schemes, Lemma
\ref{perfect-lemma-quasi-coherence-direct-image})
and is supported on $Y \setminus V$
(Derived Categories of Schemes, Definition
\ref{perfect-definition-supported-on}).
Applying $a$ we obtain a distinguished triangle
$$
a(K) \to a(Rj_*K|_V) \to a(Q) \to a(K)[1]
$$
on $X$. If $a(Q)$ is supported on $X \setminus U$, then
restricting to $U$ the map $a(K)|_U \to a(Rj_*K|_V)|_U$ is an
isomorphism, i.e., (\ref{equation-sheafy}) is an isomorphism on $K$.
The converse is immediate.

\medskip\noindent
The proof of (2) is exactly the same as the proof of (1).

\medskip\noindent
Proof of (3). Assume the equivalent conditions of (1) hold.
Set $T = Y \setminus V$.
We will use the notation $D_{\QCoh, T}(\mathcal{O}_Y)$ and
$D_{\QCoh, f^{-1}(T)}(\mathcal{O}_X)$ to denote complexes
whose cohomology sheaves are supported on $T$ and $f^{-1}(T)$.
Since $a$ commutes with direct sums, the strictly full, saturated, triangulated
subcategory $\mathcal{D}$ with objects
$$
\{Q \in D_{\QCoh, T}(\mathcal{O}_Y) \mid
a(Q) \in D_{\QCoh, f^{-1}(T)}(\mathcal{O}_X)\}
$$
is preserved by direct sums and hence derived colimits.
On the other hand, the category $D_{\QCoh, T}(\mathcal{O}_Y)$
is generated by a perfect object $E$
(see Derived Categories of Schemes, Lemma
\ref{perfect-lemma-generator-with-support}).
By assumption we see that $E \in \mathcal{D}$.
By Derived Categories, Lemma \ref{derived-lemma-write-as-colimit}
every object $Q$ of $D_{\QCoh, T}(\mathcal{O}_Y)$ is a derived
colimit of a system $Q_1 \to Q_2 \to Q_3 \to \ldots$
such that the cones of the transition maps are direct sums
of shifts of $E$. Arguing by induction we see that
$Q_n \in \mathcal{D}$ for all $n$ and finally that $Q$ is
in $\mathcal{D}$. Thus the equivalent conditions of (2) hold.
\end{proof}

\end{document}
